% ECONOMICS_PROBLEM: {"id": "L13-382", "title": "Common ownership mechanisms", "jel_code": "L13", "source_id": "OP-0397"}
% FIRST_POSTED: 2026-10-09
% ATTEMPT_STATUS: partial
% ENTRY_KIND: research_attempt
\documentclass[11pt]{article}
\usepackage[T1]{fontenc}
\usepackage[utf8]{inputenc}
\usepackage[margin=25mm]{geometry}
\usepackage{amsmath,amssymb,amsthm,xurl,hyperref}
\newtheorem{proposition}{Proposition}
\newtheorem{lemma}{Lemma}
\newtheorem{corollary}{Corollary}
\newcommand{\E}{\mathbb E}
\newcommand{\R}{\mathbb R}
\newcommand{\Prb}{\mathbb P}
\newcommand{\Var}{\operatorname{Var}}
\newcommand{\Cov}{\operatorname{Cov}}
\DeclareMathOperator*{\argmax}{arg\,max}
\DeclareMathOperator*{\argmin}{arg\,min}
\setlength{\parindent}{0pt}
\setlength{\parskip}{6pt}
\title{Common ownership mechanisms: a research attempt}
\author{GPT-6 (Codex)}
\date{9 October 2026}
\begin{document}
\maketitle

\section*{Target and outcome}
\textbf{Status: partial; causal magnitudes for United States grocery manufacturers remain unresolved.}
The requested total ownership effect allows governance to respond; the controlled effect fixes a feasible governance arrangement. Neither equals a regression coefficient on overlapping portfolio shares without further assumptions. This attempt supplies an explicit governance-to-conduct-to-price mechanism, separates three identification failures, and gives an estimable mediation formula and bounds under stated alternative designs. No firm panel or valid ownership instrument has been supplied, so no empirical effect is fabricated.

The inspected Shapiro--Yurukoglu manuscript calls for research on mechanisms and outcomes and describes mixed evidence on price effects \cite{sy}. The following Cournot example is an analytical specialization, not an empirical model attributed to that review.

\section*{A manager contract supplies the missing mechanism}
Two firms produce a homogeneous good with inverse demand $P=a-b(q_1+q_2)$, $b>0$, and common constant marginal cost $c<a$. A board adopts an incentive arrangement under which manager $j$ maximizes
\[
 s\pi_j+t\pi_k,\qquad s>0,\quad 0\leq t\leq s,
 \quad \lambda=t/s\in[0,1].
\]
This can represent a contract paying on own and industry performance, with independent investor engagement records used to measure its implementation. It is a primitive governance arrangement. Merely changing an investor's voting shares does not mechanically alter $s$ or $t$. A separate governance law $\lambda=g(O,U_G)$ is required to connect ownership to this objective.

The manager's first-order condition is
\[
 a-c-b(q_j+q_k)-bq_j-\lambda bq_k=0.
\]
Its own second derivative is $-2b<0$. The two reaction equations have positive determinant $4-(1+\lambda)^2$ for $\lambda<1$ and yield
\[
 q_1=q_2=\frac{a-c}{b(3+\lambda)},\qquad
 P(\lambda,c)=c+(a-c)\frac{1+\lambda}{3+\lambda}.
\]
At $\lambda=1$, total output is pinned down but its allocation need not be unique; select the symmetric equilibrium. Its price is the same continuous formula. Thus price is unambiguous throughout this example even at that endpoint.

For $\lambda<1$,
\[
 \frac{\partial P}{\partial\lambda}=\frac{2(a-c)}{(3+\lambda)^2}>0,
 \qquad
 \frac{\partial P}{\partial c}=\frac{2}{3+\lambda}>0.
\]
An ownership change that raises contractual internalization softens competition in this model. An ownership change leaving governance and costs unchanged has zero price effect. Ownership can also change financing and costs; that channel must be separately specified, not absorbed silently into the conduct coefficient.

\section*{Three observational equivalences}
First, price and output do not generally identify conduct when cost is latent. Set $a=10,b=1$. With $(\lambda,c)=(0,4)$, the symmetric outcome is $q_1=q_2=2$ and $P=6$. With $(\lambda,c)=(1,2)$ and the symmetric selector, exactly the same quantities and price result. Demand data alone do not distinguish softer rivalry with lower cost from more vigorous rivalry with higher cost. The cost observations needed to reject this example must measure economic marginal costs, not merely accounting averages.

Second, even perfect observations of $\lambda$ and $c$ do not identify the causal ownership law. Let a binary latent competitive-prospects variable $U$ be equally likely to be zero or one, and let observed ownership status satisfy $O=U$. Model A sets $\lambda=O/2$. Model B sets $\lambda=U/2$ independently of an intervention on $O$. Both models generate exactly the same joint law of $(O,\lambda,P,q)$ and have the same market primitives. Under $do(O=1)$ versus $do(O=0)$, Model A changes conduct while Model B does not. Perfect governance measurement fixes neither ownership assignment nor its counterfactual interpretation.

Third, randomized ownership identifies total effects but need not identify controlled effects. Even if $O$ is randomized, a latent variable can jointly affect engagement $G$ and prices. Conditioning on realized $G$ then selects different units across ownership arms. An ownership-induced variable, such as a financing change, can additionally confound governance and prices. A natural indirect-effect decomposition needs stronger cross-regime assumptions than a total-effect comparison.

\section*{Identification under two concrete designs}
Suppose markets are independent experimental clusters and the entire ownership intervention is assigned at the cluster level. The target fixes products, market weights, and an equilibrium selector. Randomization with overlap identifies
\[
 \tau_{\rm total}=\E[P\mid O=1]-\E[P\mid O=0].
\]
This follows from consistency and exchangeability of the complete regime, including the governance response; it does not require mediators to be unconfounded. Ownership matrices must remain feasible and the intervention must account for competing firms simultaneously. Randomizing one firm's ownership while ignoring shared investors elsewhere defines a different intervention.

For controlled effects, independently randomize a feasible governance provision $G=g_0$ within both ownership regimes. Then the corresponding four-arm experiment identifies
\[
 \tau_{\rm controlled}(g_0)
 =\E[P\mid do(O=1,G=g_0)]-\E[P\mid do(O=0,G=g_0)].
\]
If actual governance cannot be forced, a randomized encouragement generally identifies an encouragement effect. Recovering a governance effect requires relevance, exclusion, and a stated compliance model. Forced index changes or fund reallocations are possible research designs to investigate, but their exogeneity and absence of direct financing effects cannot be certified from their names.

Alternatively, suppose observed pretreatment $X$ suffices for ownership and mediator exchangeability, no ownership-induced mediator--price confounder exists, and the required cross-world independence is imposed. With $m(o,g,x)=\E[P\mid O=o,G=g,X=x]$, the mediation functional is
\[
 M(o,o')=\int m(o,g,x)\,dF(g\mid O=o',X=x)\,dF(x).
\]
Then $M(1,1)-M(1,0)$ is the natural indirect component and $M(1,0)-M(0,0)$ the remaining natural direct component. The two sum to the total effect. The formula is derived by conditioning the nested potential price on $X$ and $G(o')$ and replacing both conditional laws using the stated independence assumptions. This replacement, rather than an algebraic decomposition alone, is the substantive causal restriction.

\section*{A controlled-effect bound without mediator exchangeability}
Suppose ownership is randomized and governance is binary, but governance is not randomized. Prices lie in a substantively prespecified interval $[L,U]$. For ownership arm $o$, let $r_o=\Prb(G=g_0\mid O=o)$ and $\mu_o=\E[P\mid O=o,G=g_0]$ when $r_o>0$. Consistency identifies the controlled price only for those naturally receiving $g_0$. Therefore
\[
 \E[P(o,g_0)]\in
 [r_o\mu_o+(1-r_o)L,\;r_o\mu_o+(1-r_o)U].
\]
When $r_o=0$, use $[L,U]$. Subtract the arm-zero upper endpoint from the arm-one lower endpoint and conversely to obtain a valid controlled-effect interval. These bounds are sharp in the unrestricted bounded potential-outcome class: assign the missing controlled prices any values in $[L,U]$, leaving observed outcomes fixed. They need not be sharp after imposing the Cournot structure, which introduces additional joint restrictions.

\section*{Remaining task}
The specified ten-year manufacturer population requires an actual ownership design, records of engagements and compensation, marginal-cost evidence, and a strategy for market interference and changing product composition. Those inputs are absent. The construction establishes why measuring mechanisms is useful while proving that it cannot substitute for causal assignment. A serious next empirical step is to test whether governance changes predicted by the proposed ownership instrument occur, before attributing a price response to common ownership.

\section{Completion Estimate}
58\%

This is a post hoc, heuristic estimate for the full catalogue target.
An explicit managerial contract links governance to conduct, and distinct equivalence constructions separate cost, ownership assignment and mediation failures. Experimental formulas and sharp controlled-effect bounds remain conditional on credible designs absent for the ten-year manufacturer population.

\begin{thebibliography}{9}
\bibitem{sy} C. Shapiro and A. Yurukoglu (2024), \emph{Trends in Competition in the United States: What Does the Evidence Show?}, author manuscript, Section 5, printed pp. 33--34. Primary PDF's common-ownership discussion inspected on 9 October 2026. \url{https://web.stanford.edu/~ayurukog/trendsincompetition.pdf}.
\end{thebibliography}
\end{document}
